\section{Alcance y criterio de evidencia}

El análisis distingue cuatro niveles de evidencia:

\begin{enumerate}
  \item \textbf{derivación}: existencia de una rama armónica compatible y de
        una semilla obtenida desde el modelo;
  \item \textbf{localización}: integración y continuación hasta un conjunto
        recurrente del sistema original;
  \item \textbf{clasificación dinámica}: periodicidad, regularidad o
        comportamiento compatible con caos bajo un diagnóstico finito;
  \item \textbf{ocultedad}: ausencia de contacto desde un conjunto explícito
        de vecindades de todos los equilibrios.
\end{enumerate}

La continuación establece la localización de un candidato. La clasificación
dinámica requiere diagnósticos convergentes y la compatibilidad con ocultedad
queda limitada al banco de vecindades y al horizonte declarados.

\evidencebox{\textbf{Interpretación.} \emph{Compatible con ocultedad bajo
el protocolo finito} significa exactamente que ninguna condición inicial del
banco ensayado alcanzó la nube candidata. \emph{Oculto según la fuente}
atribuye el rótulo a la bibliografía. \emph{Ocultedad certificada en un
dominio} se reserva para un argumento que cubra el continuo declarado. Ninguna
de las dos primeras expresiones significa una prueba topológica global de la
cuenca.}

\section{Ruta matemática y numérica común}

\subsection{Declaración Lur'e y ruta directa D0}

La formulación escalar base es
\begin{equation}
  \dot X=PX+b\,\psi(\sigma),\qquad
  \sigma=c^{\mathsf T}X,
  \label{int:eq:lure}
\end{equation}
con una comprobación numérica de que el campo de \eqref{int:eq:lure} coincide con
las ecuaciones originales. Se verifican también los equilibrios, el Jacobiano
analítico contra diferencias finitas y el espectro local.

La transferencia se define por
\begin{equation}
  G(s)=c^{\mathsf T}(P-sI)^{-1}b.
  \label{int:eq:transfer}
\end{equation}
Para justificar el signo del cierre, considérese primero la ganancia lineal
\(\psi(\sigma)=k\sigma\) y un modo \(X(t)=v\e^{st}\). Entonces
\[
 (sI-P)v=bk\,c^{\mathsf T}v.
\]
Si \(sI-P\) es invertible y \(c^{\mathsf T}v\ne0\), multiplicar por
\(c^{\mathsf T}(sI-P)^{-1}\) y cancelar \(c^{\mathsf T}v\) da
\[
 1=k\,c^{\mathsf T}(sI-P)^{-1}b=-kG(s).
\]
Equivalentemente, el lema del determinante para una perturbación de rango uno
da
\[
 \det(sI-P-kbc^{\mathsf T})=\det(sI-P)[1+kG(s)].
\]
Estas identidades se aplican fuera de los polos de la transferencia. Para
\(s=i\omega_0\) y \(k\) real, separarlas en partes real e imaginaria
produce el cierre siguiente. La ganancia de primer armónico de una entrada
\(a\cos\theta\) es
\[
 \mathcal N(a)=\frac1{\pi a}\int_0^{2\pi}
                   \psi(a\cos\theta)\cos\theta\,\dd\theta.
\]
Se obtiene por proyección sobre \(\cos\theta\), cuya norma cuadrática en
un periodo es \(\pi\). El cierre retiene esa componente; los armónicos
superiores se evalúan posteriormente mediante la integración del campo.

La ruta directa resuelve
\begin{equation}
  \operatorname{Im}G(i\omega_0)=0,\qquad
  k=-\frac{1}{\operatorname{Re}G(i\omega_0)},\qquad
  \mathcal N(a_0)=k.
  \label{int:eq:balance}
\end{equation}
La frecuencia \(\omega_0\) se obtiene de las raíces polinómicas o de la
ecuación escalar derivada de la transferencia. El diagrama de Nyquist
representa las ramas resultantes.

Para cada rama compatible se forma \(P_0=P+kbc^{\mathsf T}\). Un vector
normalizado se obtiene directamente como
\[
 v=k(i\omega_0I-P)^{-1}b,\qquad
 c^{\mathsf T}v=-kG(i\omega_0)=1.
\]
La ecuación resolvente implica \(P_0v=i\omega_0v\). Por ser \(P_0\) real,
\[
 X_{\rm aux}(t)=a_0\operatorname{Re}
       \bigl(v\e^{i(\omega_0t+\theta)}\bigr),\qquad
 c^{\mathsf T}X_{\rm aux}(t)=a_0\cos(\omega_0t+\theta)
\]
es una solución exacta del sistema auxiliar. Su estado inicial es la semilla
\(X_0=a_0\operatorname{Re}(v\e^{i\theta})\). La continuación entera es
\begin{equation}
  \dot X=P_0X+\lambda b
  \left[\psi(c^{\mathsf T}X)-k\,c^{\mathsf T}X\right],
  \qquad 0\leq\lambda\leq1.
  \label{int:eq:homotopy}
\end{equation}
En $\lambda=0$ queda el oscilador armónico auxiliar y en $\lambda=1$ se
recupera exactamente el sistema original.

\subsection{Ruta alternativa A1}

Si \eqref{int:eq:balance} no produce una ganancia compatible, una amplitud real,
una continuación persistente o candidatos que superen el objetivo dinámico
declarado, se recurre a una alternativa adecuada a esa obstrucción. En Kalman--Fitts, la única rama directa tiene
ganancia negativa, incompatible con la función descriptiva positiva de
\(\tanh\). Entonces se usa un mapa de conmutación de Andronov para
\(\operatorname{sign}(\sigma)\) y una homotopía
\begin{equation}
  \psi_\lambda(\sigma)=
  (1-\lambda)\operatorname{sign}(\sigma)
  +\lambda\tanh(\sigma/\varepsilon).
  \label{int:eq:sign-tanh}
\end{equation}

Los barridos de frecuencia, las transferencias sesgadas y la continuación
multiparámetro son alternativas reservadas para fallos explícitos. No forman
parte de la ruta principal de orden entero. En la extensión caótica MAVPD, la
ruta D0 sí produce semillas y continuación, pero ninguna de sus dos ramas base
supera el cribado finito de caos; sólo entonces se habilita una continuación
paramétrica estructurada alrededor de una frontera de Hopf derivada.

\subsection{Localización y contraste independiente}

Las semillas se calculan a partir de las ecuaciones del modelo. Los valores
publicados y los obtenidos previamente mediante cálculo simbólico se emplean
para contrastar el cierre algebraico y el destino de las trayectorias una
vez completada la localización. La comparación entre Python y Julia utiliza
cálculos independientes con un protocolo común de evaluación.

\subsection{Protocolo de ocultedad}

Después de localizar una nube candidata \(\mathcal A_h\), se construyen
condiciones iniciales alrededor de \emph{todos} los equilibrios. Para una
condición \(X_0\), la cola de su trayectoria se compara con:

\begin{enumerate}
  \item las tolerancias de convergencia a cada equilibrio;
  \item la nube candidata mediante una distancia normalizada;
  \item umbrales de divergencia y de conjunto acotado no trivial.
\end{enumerate}

La conclusión \emph{sin contacto} queda ligada al conjunto explícito de
condiciones iniciales, junto con sus radios, direcciones, paso, transitorio,
horizonte y tolerancia geométrica. En cada comparación cruzada los dos entornos
reciben las mismas coordenadas.

\begin{center}
\begin{tikzpicture}[
  node distance=0.55cm and 0.9cm,
  box/.style={draw=Azul,rounded corners,fill=AzulClaro,align=center,
              text width=3.45cm,minimum height=0.88cm},
  alt/.style={draw=Naranja,rounded corners,fill=NaranjaClaro,align=center,
              text width=3.45cm,minimum height=0.88cm},
  arrow/.style={-{Latex[length=2.3mm]},thick}
]
\node[box] (model) {Ecuaciones, Jacobiano,\\equilibrios y forma Lur'e};
\node[box,below=of model] (direct) {Transferencia directa\\$\omega_0\to k\to a_0\to X_0$};
\node[box,below=of direct] (cont) {Continuación al\\sistema original};
\node[box,below=of cont] (dyn) {Candidato y\\diagnóstico dinámico};
\node[box,below=of dyn] (hidden) {Sondeos comunes\\de ocultedad};
\node[alt,right=of direct] (fallback) {Ruta alternativa sólo\\tras fallo documentado};
\draw[arrow] (model)--(direct);
\draw[arrow] (direct)--(cont);
\draw[arrow] (cont)--(dyn);
\draw[arrow] (dyn)--(hidden);
\draw[arrow,dashed,draw=Naranja] (direct)--(fallback);
\draw[arrow,dashed,draw=Naranja] (fallback.south)--(dyn.east);
\end{tikzpicture}

\smallskip
\small\textbf{Diagrama metodológico.} Jerarquía de búsqueda usada para los
sistemas enteros.
\end{center}

\section[Caso 1: Chua no suave entero]{Caso 1: Chua no suave de orden entero}

\subsection{Modelo y derivación}

El sistema es
\begin{align}
 \dot x&=\alpha\,[y-x-\phi(x)],\nonumber\\
 \dot y&=x-y+z,\nonumber\\
 \dot z&=-\beta y-\gamma z, \label{int:eq:chua}\\
 \phi(x)&=m_1x+(m_0-m_1)\operatorname{sat}(x),\qquad
 \operatorname{sat}(x)=\min(1,\max(-1,x)).\nonumber
\end{align}
Se usaron
\[
 \alpha=8.4562,\quad \beta=12.0732,\quad \gamma=0.0052,\quad
 m_0=-0.1768,\quad m_1=-1.1468.
\]
La forma D0 exacta toma \(c=e_1\), \(b=-\alpha e_1\), una parte lineal
con pendiente base \(m_1\) y
\(\psi(\sigma)=(m_0-m_1)\operatorname{sat}(\sigma)\).

Escribiendo \(a_P=\alpha(1+m_1)\), la matriz lineal y el polinomio
intermedio son
\[
 P=\begin{pmatrix}-a_P&\alpha&0\\1&-1&1\\0&-\beta&-\gamma\end{pmatrix},
 \qquad D(s)=s^2+(1+\gamma)s+\beta+\gamma.
\]
Las dos últimas filas de \((sI-P)v=b\) dan
\(v_3=-\beta v_2/(s+\gamma)\) y
\(v_1=D(s)v_2/(s+\gamma)\). La primera fila produce
\[
 [(s+a_P)D(s)-\alpha(s+\gamma)]v_1=-\alpha D(s).
\]
Por tanto,
\begin{equation}
 G(s)=\frac{\alpha D(s)}{(s+a_P)D(s)-\alpha(s+\gamma)}.
 \label{int:eq:chua-transfer-derived}
\end{equation}
Con \(d_0=\beta+\gamma\), \(d_1=1+\gamma\), \(z=\omega^2\), el numerador
complejo de \(G(i\omega)\) es \(\alpha(d_0-z+id_1\omega)\), y el
denominador tiene partes
\[
 A=a_Pd_0-\alpha\gamma-(d_1+a_P)z,\qquad
 B=\omega(d_0+a_Pd_1-\alpha-z).
\]
Como \(\operatorname{Im}G=\alpha[d_1\omega A-(d_0-z)B]/(A^2+B^2)\),
la condición imaginaria, para \(\omega>0\) y \(A^2+B^2>0\), equivale a
\[
 z^2+(\alpha+d_1^2-2d_0)z+d_0^2-\alpha d_0+\alpha\gamma d_1=0.
\]
Cada raíz positiva se comprueba en el resolvente antes de calcular \(k\).
Normalizando el primer componente del modo a uno, sus otras componentes son
\[
 v_2=\frac{i\omega+\gamma}{D(i\omega)},\qquad
 v_3=-\frac{\beta}{D(i\omega)}.
\]
Así se obtienen los tres componentes de la semilla a partir de la amplitud y
la fase, sin ajustar coordenadas independientemente.

Sea \(\Delta=m_0-m_1\). Para \(0<a\leq1\), la saturación permanece en su
tramo lineal y \(\mathcal N(a)=\Delta\). Para \(a>1\), el ángulo de cruce
es \(\theta_c=\arccos(1/a)\). La simetría impar y la simetría entre
cuadrantes reducen la integral a
\begin{align*}
 \mathcal N(a)
 &=\frac{4\Delta}{\pi a}
   \left[\int_0^{\theta_c}\cos\theta\,\dd\theta
       +a\int_{\theta_c}^{\pi/2}\cos^2\theta\,\dd\theta\right]\\
 &=\frac{4\Delta}{\pi a}
   \left[\sin\theta_c+a\left(\frac\pi4-\frac{\theta_c}{2}
                              -\frac{\sin2\theta_c}{4}\right)\right].
\end{align*}
Usando \(\pi/2-\theta_c=\arcsin(1/a)\),
\(\sin\theta_c=\sqrt{a^2-1}/a\) y
\(\sin2\theta_c=2\sqrt{a^2-1}/a^2\), se obtiene
\begin{equation}
 \mathcal N(a)=\frac{2(m_0-m_1)}{\pi}
 \left[
 \arcsin\!\left(\frac1a\right)
 +\frac{\sqrt{a^2-1}}{a^2}
 \right].
\end{equation}
El balance armónico produjo
\[
 \omega_0=2.039186939959,\qquad
 k=0.209867354515,\qquad
 a_0=5.856145086252,
\]
y la semilla
\[
 X_0=(5.856145086252,\ 0.369331578246,\ -8.366536168324)^{\mathsf T}.
\]
El residuo de cierre de Nyquist fue \(3.26\times10^{-15}\) y el de la
función descriptiva \(1.98\times10^{-13}\). El contraste posterior con los cálculos de
Mathematica y MATLAB dio diferencias máximas del orden de \(10^{-12}\).

\subsection{Localización por recurrencias}

La búsqueda por recurrencias en Julia evaluó 512 condiciones iniciales en
un dominio y una partición prefijados. Encontró dos conjuntos recurrentes,
identificados mediante las etiquetas 2 y 3, con 22 contactos entre los 512
lanzamientos: una fracción muestral de \(0.04296875\). Las semillas de la
función descriptiva confirmaron posteriormente la correspondencia con las
ramas positiva y negativa.

\begin{hafocommonfigure}[H]
\centering
\begin{subfigure}[t]{0.49\textwidth}
  \centering
  \includegraphics[width=\linewidth]{chua_python_attractor.png}
  \caption{\hafo{}: candidato obtenido por semilla y continuación.}
\end{subfigure}\hfill
\begin{subfigure}[t]{0.49\textwidth}
  \centering
  \includegraphics[width=\linewidth]{chua_julia_candidates.png}
  \caption{Julia: dos conjuntos descubiertos por recurrencias globales.}
\end{subfigure}
\caption{Localización independiente en el Chua no suave.}
\label{int:fig:chua-localization}
\end{hafocommonfigure}

\subsection{Comparación dinámica Python--Julia}

\ifHAFOBasinMaterial
\paragraph{Protocolo de comparación cruzada.}
Ambos cálculos evaluaron \ChuaProbeCount{} condiciones:
3 equilibrios, 7 radios entre \(10^{-5}\) y \(10^{-2}\), y 24 puntos por
equilibrio y radio. Se usaron \(t_f=500\), \(t_{\mathrm{burn}}=120\),
tolerancia de equilibrio \(10^{-3}\) y tolerancia normalizada de nube \(0.08\).
Las 504 coordenadas se generaron uniformemente en bolas, con semilla
pseudoaleatoria 123456789. El control EFORK--3 de referencia emplea
\(T=320\), \(T_{\mathrm{burn}}=80\), otras coordenadas y un clasificador
por huellas de sección.
\fi

\begin{table}[H]
\centering
\caption{Comparación Python--Julia del caso de Chua.}
\small
\begin{tabularx}{\textwidth}{@{}lYY@{}}
\toprule
Magnitud & Python/\hafo{} & Julia \\
\midrule
\ifHAFOBasinMaterial
Clases finales coincidentes & \multicolumn{2}{c}{504 de 504} \\
Conteos comunes & \multicolumn{2}{c}{152 en \(E_0\), 198 acotadas no triviales, 154 divergentes} \\
Contactos con candidato & 0 & 0 \\
\fi
Exponentes de Lyapunov finitos &
\((0.14018,\ 0.01319,\ -1.20304)\) &
\((0.14591,\ -0.00052,\ -1.17115)\) \\
\bottomrule
\end{tabularx}
\end{table}

\ifHAFOBasinMaterial
Los conteos \(152/198/154\) pertenecen exclusivamente a este protocolo
de comparación cruzada. Los conteos \(369/69/66\) de la configuración de referencia usan
otro horizonte, otras coordenadas y otro clasificador. Ambos protocolos dieron cero contactos con la nube objetivo.
\fi

Los dos espectros son compatibles con un exponente máximo positivo. Esa
observación apoya la clasificación caótica local, pero no sustituye la prueba
de ocultedad.

\begin{hafobasinfigure}[H]
\centering
\includegraphics[width=0.82\textwidth]{chua_python_hiddenness.png}
\caption{Control de vecindades de los tres equilibrios en Python. El
comparador usó las mismas 504 coordenadas en Julia.}
\label{int:fig:chua-hiddenness}
\end{hafobasinfigure}

\section[Caso 2: Kalman--Fitts no Chua]
{Caso 2: Kalman--Fitts, sistema no Chua de cuarto orden}

\subsection{Formulación y rechazo de la ruta directa}

El sistema exacto es
\begin{equation}
 \dot X=AX+b\tanh\!\left(\frac{c^{\mathsf T}X}{\varepsilon}\right),
 \quad b=e_4,\quad c=-e_3,\quad \varepsilon=0.01,
 \label{int:eq:kf}
\end{equation}
La matriz compañera se define por
\[
 A=\begin{pmatrix}
 0&1&0&0\\0&0&1&0\\0&0&0&1\\-a_0&-a_1&-a_2&-a_3
 \end{pmatrix},
 \qquad
 p_A(s)=\bigl[(s+\beta)^2+m_1^2\bigr]
         \bigl[(s+\beta)^2+m_2^2\bigr].
\]
Al expandir el producto,
\begin{align*}
 a_0&=(m_1^2+\beta^2)(m_2^2+\beta^2),&
 a_1&=2\beta(m_1^2+m_2^2+2\beta^2),\\
 a_2&=m_1^2+m_2^2+6\beta^2,&a_3&=4\beta.
\end{align*}
Los parámetros publicados \(m_1=0.9\), \(m_2=1.1\), \(\beta=0.03\)
\cite{kalman2019} dan
\((a_0,a_1,a_2,a_3)=(0.98191881,0.121308,2.0254,0.12)\).
Las tres primeras ecuaciones de equilibrio imponen
\(x_2=x_3=x_4=0\); la cuarta queda \(-a_0x_1=0\), de modo que el único
equilibrio es el origen. Su Jacobiano es \(A+bc^{\mathsf T}/\varepsilon\)
y tiene polinomio
\[
 s^4+a_3s^3+(a_2+1/\varepsilon)s^2+a_1s+a_0.
\]
Todos los coeficientes son positivos. Las dos desigualdades adicionales de
Routh--Hurwitz para este polinomio mónico son
\[
 a_3(a_2+1/\varepsilon)-a_1=12.12174>0,
\]
\[
 a_3(a_2+1/\varepsilon)a_1-a_1^2-a_3^2a_0
 =1.456324405056>0.
\]
Por tanto, el equilibrio es localmente asintóticamente estable.

Para calcular la transferencia se resuelve \((sI-A)v=b\): las tres primeras
filas dan \(v_2=sv_1\), \(v_3=s^2v_1\), \(v_4=s^3v_1\), y la última da
\(p_A(s)v_1=1\). Como \(c=-e_3\),
\[
 c^{\mathsf T}(sI-A)^{-1}b=-\frac{s^2}{p_A(s)},\qquad
 G(s)=\frac{s^2}{p_A(s)}.
\]
En \(s=i\omega\), las partes real e imaginaria del denominador son
\(R=\omega^4-a_2\omega^2+a_0\) e
\(I=\omega(a_1-a_3\omega^2)\). De aquí
\[
 \operatorname{Im}G(i\omega)=\frac{\omega^2I}{R^2+I^2},\qquad
 \omega_0^2=\frac{a_1}{a_3},\qquad
 k=\frac{\omega_0^4-a_2\omega_0^2+a_0}{\omega_0^2}.
\]
La ruta directa, sin malla de frecuencias, encontró
\[
 \omega_0=1.005435229142,\qquad k=-0.043168701157.
\]
Para la no linealidad hiperbólica,
\[
 \mathcal N(a)=\frac1{\pi\varepsilon}\int_0^{2\pi}
  \frac{\tanh u}{u}\cos^2\theta\,\dd\theta,
 \qquad u=\frac{a\cos\theta}{\varepsilon},
\]
con valor continuo \(\tanh u/u=1\) en \(u=0\). Para \(u\ne0\),
\(0<\tanh u/u<1\), de modo que \(0<\mathcal N(a)<1/\varepsilon\).
Además, \(\mathcal N(a)\to1/\varepsilon\) cuando \(a\to0^+\), mientras
\(\mathcal N(a)\leq4/(\pi a)\to0\) cuando \(a\to\infty\). La continuidad
y estas cotas dan el rango \((0,100)\).
Por tanto, la ganancia negativa se rechaza matemáticamente y no
existe amplitud directa compatible. Python y Julia coincidieron en
\(\omega_0\) a \(1.02\times10^{-14}\) y en \(k\) a
\(1.09\times10^{-15}\).

\subsection{Mapa de conmutación y continuación A1}

En cada semiespacio \(\operatorname{sign}(c^{\mathsf T}X)=\sigma\),
\(\sigma\in\{-1,1\}\), el precursor satisface \(\dot X=AX+b\sigma\).
La variación de constantes y la invertibilidad de \(A\) dan
\begin{equation}
 \Phi_\sigma(t,X)=\e^{At}X+
                   A^{-1}(\e^{At}-I)b\sigma.
 \label{int:eq:kf-switching-flow}
\end{equation}
En efecto, \(\int_0^t\e^{A(t-s)}\,\dd s=A^{-1}(\e^{At}-I)\). En la
sección \(\Sigma=\{X:c^{\mathsf T}X=0\}\), el lado de salida se elige
según el signo de \(c^{\mathsf T}(AX+b\sigma)\). El tiempo \(t_+(X)>0\)
es el primer cero transversal de
\(c^{\mathsf T}\Phi_\sigma(t,X)\) después de abandonar la sección. Así se
define el mapa de conmutación
\(\Pi(X)=\Phi_\sigma(t_+(X),X)\). Una órbita de periodo dos cumple
\(\Pi^2(X)=X\), con \(\Pi(X)\ne X\); el periodo continuo suma los dos
tiempos de vuelo. La transversalidad excluye los contactos tangenciales en
los que el siguiente cruce no queda definido de este modo.

El mapa de sección de la no linealidad \(\operatorname{sign}(\sigma)\)
convergió en 135 iteraciones a una órbita de retorno de periodo 2 y generó
\[
 X_{\mathrm{sw}}=
 (0.212347571803,\ -1.671123727247,\ 0,\ 1.478087383043)^{\mathsf T}.
\]
Después se aplicaron 51 valores de \(\lambda\) en
\eqref{int:eq:sign-tanh}. Ninguna implementación leyó el punto publicado como
semilla. El punto de la Tabla 2 se comparó sólo después y quedó a distancia
\(4.34\times10^{-4}\) de la trayectoria Python y
\(4.30\times10^{-4}\) de la Julia.

\begin{hafocommonfigure}[H]
\centering
\begin{subfigure}[t]{0.49\textwidth}
  \centering
  \includegraphics[width=\linewidth]{kalman_continuation.png}
  \caption{Homotopía \(\operatorname{sign}\to\tanh\) en Python.}
\end{subfigure}\hfill
\begin{subfigure}[t]{0.49\textwidth}
  \centering
  \includegraphics[width=\linewidth]{kalman_python_julia_overlay.png}
  \caption{Nubes finales generadas independientemente.}
\end{subfigure}
\caption{Ruta alternativa reproducida para Kalman--Fitts.}
\label{int:fig:kf}
\end{hafocommonfigure}

\subsection{Ciclo límite y comparación dinámica}

El candidato es un ciclo límite estable. El periodo medio fue
\(11.9169960\) en Python y \(11.9169380\) en Julia, con diferencia
\(5.81\times10^{-5}\). La distancia normalizada entre nubes fue
\(2.10\times10^{-3}\).

\ifHAFOBasinMaterial
El protocolo compartido contiene \KFProbeCount{} condiciones alrededor del
único equilibrio: 4 radios de \(10^{-5}\) a \(10^{-2}\) y 12 condiciones por
radio. Los dos códigos produjeron 40 convergencias al equilibrio, 8 estados
acotados no triviales y cero contactos con el ciclo. Las
\KFProbeCount{} clases finales y las \KFProbeCount{} decisiones de contacto
coincidieron punto por punto.
Por tanto, el ciclo es compatible con ocultedad bajo estas 48 condiciones
iniciales. El muestreo no prueba que la cuenca esté separada de toda vecindad
abierta del equilibrio.
\fi

\begin{table}[H]
\centering
\caption{Comparación cerrada de Kalman--Fitts.}
\small
\begin{tabularx}{\textwidth}{@{}lYY@{}}
\toprule
Magnitud & Python/\hafo{} & Julia \\
\midrule
Ruta directa & Rechazo por \(k<0\) & Rechazo por \(k<0\) \\
Mapa de conmutación & 135 iteraciones, retorno 2 & 135 iteraciones, retorno 2 \\
Periodo del ciclo & 11.9169960 & 11.9169380 \\
\ifHAFOBasinMaterial
Contactos en 48 sondeos & 0 & 0 \\
\fi
\bottomrule
\end{tabularx}
\end{table}


\section[Caso 3: Van der Pol--Duffing modificado]
{Caso 3: Van der Pol--Duffing autónomo modificado}

\subsection{Modelo y forma D0 exacta}

El sistema MAVPD de orden entero es
\begin{align}
 \dot y_1&=\delta\gamma y_1+\delta y_2-\delta y_1^3,\nonumber\\
 \dot y_2&=y_1-\xi y_2-y_3,\label{int:eq:mavpd}\\
 \dot y_3&=\rho y_2,\nonumber
\end{align}
con \(\gamma=0.1\), \(\delta=100\) y \(\rho=200\)
\cite{matouk2023}. Su declaración escalar es
\[
 P=\begin{pmatrix}
 \delta\gamma&\delta&0\\
 1&-\xi&-1\\
 0&\rho&0
 \end{pmatrix},\quad
 b=\begin{pmatrix}-\delta\\0\\0\end{pmatrix},\quad
 c=\begin{pmatrix}1\\0\\0\end{pmatrix},\quad
 \psi(\sigma)=\sigma^3,
\]
de modo que el Jacobiano y todos los equilibrios reales para \(\gamma>0\) son
\begin{equation}
 J(y)=\begin{pmatrix}
 \delta(\gamma-3y_1^2)&\delta&0\\
 1&-\xi&-1\\
 0&\rho&0
 \end{pmatrix},\qquad
 E_0=(0,0,0),\quad
 E_\pm=(\pm\sqrt{\gamma},0,\pm\sqrt{\gamma}).
 \label{int:eq:mavpd-jacobian}
\end{equation}
Para obtener los equilibrios, \(\dot y_3=0\) exige \(y_2=0\) porque
\(\rho>0\); \(\dot y_2=0\) da \(y_3=y_1\); la primera ecuación se reduce
a \(\delta y_1(\gamma-y_1^2)=0\). Esto produce exactamente los tres puntos
anteriores cuando \(\gamma>0\).

Sea \(u=(P-sI)^{-1}b\). Las últimas filas de este sistema lineal son
\[
 u_1-(\xi+s)u_2-u_3=0,\qquad \rho u_2-su_3=0.
\]
Al eliminar \(u_3\), resulta
\(u_2=su_1/(s^2+\xi s+\rho)\). La primera fila queda
\[
 \left[\delta\gamma-s+\frac{\delta s}{s^2+\xi s+\rho}\right]u_1=-\delta.
\]
Despejando \(u_1=G(s)\) y expandiendo
\((s-\delta\gamma)(s^2+\xi s+\rho)-\delta s\), se obtiene
\begin{equation}
 G(s)=\frac{\delta(s^2+\xi s+\rho)}
 {s^3+(\xi-\delta\gamma)s^2+
 (\rho-\delta-\delta\gamma\xi)s-\delta\gamma\rho}.
 \label{int:eq:mavpd-transfer}
\end{equation}
El caso Wolfram conserva además la convención resolvente
\[
 W(s)=c^{\mathsf T}(sI-P)^{-1}b=-G(s).
\]
Por tanto, su cálculo \(k=1/\operatorname{Re}W(i\omega)\) es exactamente el
mismo que \(k=-1/\operatorname{Re}G(i\omega)\) en
\eqref{int:eq:mavpd-direct-scalars}; el signo queda fijado por la derivación
algebraica.
La identidad \(\cos^3\theta=(3\cos\theta+\cos3\theta)/4\) separa las
componentes de \(\psi(a\cos\theta)=a^3\cos^3\theta\). El primer armónico
tiene amplitud \(3a^3/4\); dividir entre la amplitud de entrada \(a\) da
\begin{equation}
  \mathcal N(a)=\frac{3a^2}{4}.
\end{equation}

Al escribir \(z=\omega^2\), el numerador sin el factor \(\delta\) tiene
partes \(\rho-z\) y \(\xi\omega\), mientras el denominador tiene partes
\[
 C=-\delta\gamma\rho-(\xi-\delta\gamma)z,\qquad
 D=\omega(\rho-\delta-\delta\gamma\xi-z).
\]
Racionalizar el cociente produce
\[
 \operatorname{Im}G(i\omega)=
 \frac{\delta[\xi\omega C-(\rho-z)D]}{C^2+D^2}.
\]
Al expandir el numerador, los términos con \(\gamma\) se cancelan y queda
\[
 \xi\omega C-(\rho-z)D
 =-\omega[z^2+(\delta-2\rho+\xi^2)z+\rho(\rho-\delta)].
\]
Para \(\omega>0\) y \(C^2+D^2>0\), la condición imaginaria equivale a
\begin{equation}
 z^2+(\delta-2\rho+\xi^2)z+\rho(\rho-\delta)=0.
 \label{int:eq:mavpd-poly}
\end{equation}
Así se obtienen todas las ramas sin barrido de frecuencias. Para cada raíz
positiva \(z_j\), el procedimiento calcula
\begin{equation}
 \omega_j=\sqrt{z_j},\qquad
 k_j=-\frac{1}{\operatorname{Re}G(i\omega_j)},\qquad
 a_j=\sqrt{\frac{4k_j}{3}},
 \label{int:eq:mavpd-direct-scalars}
\end{equation}
siempre que \(k_j>0\), condición necesaria para una amplitud real positiva.
Para el modo normalizado \(v_{j,1}=1\), la primera fila de
\((P+k_jbc^{\mathsf T})v_j=i\omega_jv_j\) es
\[
 \delta(\gamma-k_j)+\delta v_{j,2}=i\omega_j,
\]
y la tercera es \(\rho v_{j,2}=i\omega_jv_{j,3}\). Despejar ambas da
\begin{equation}
 v_j=\begin{pmatrix}
 1\\
 k_j-\gamma+i\omega_j/\delta\\
 \rho/\delta-i\rho(k_j-\gamma)/\omega_j
 \end{pmatrix},\qquad
 y_{0,j}(\theta)=a_j\operatorname{Re}\!\left(v_je^{i\theta}\right).
 \label{int:eq:mavpd-direct-seed}
\end{equation}
Las igualdades de cierre de \eqref{int:eq:mavpd-direct-scalars}, la ecuación modal
\((P+k_jbc^{\mathsf T})v_j=i\omega_jv_j\) y el residuo del campo se guardan
como verificaciones, no como datos impuestos.

\subsection{\texorpdfstring{\(\xi=3.1\): ramas, semillas y continuación}{xi=3.1: ramas, semillas y continuación}}

Para \(\xi=3.1\), \eqref{int:eq:mavpd-poly} tiene coeficientes
\((1,-290.39,20000)\) y produce dos ramas compatibles:

\begin{table}[H]
\centering
\caption{Ramas directas MAVPD para \(\xi=3.1\).}
\small
\begin{tabular}{@{}ccccc@{}}
\toprule
Rama & \(\omega_0\) & \(k\) & \(a_0\) & Semilla con fase \(0\) \\
\midrule
0 & 10.5975230311 & 0.1397015948 & 0.4315886851 &
\((0.4315887,0.0171348,0.8631774)\) \\
1 & 13.3447557342 & 0.3518790503 & 0.6849613618 &
\((0.6849614,0.1725274,1.3699227)\) \\
\bottomrule
\end{tabular}
\end{table}

Se ejecutaron ambas ramas con fases \(0\) y \(\pi\), y 21 valores
\(\lambda=0,0.05,\ldots,1\). La rama 0 conduce a dos ciclos internos
simétricos; la rama 1 conduce al conjunto externo. La selección se hizo por
mayor extensión de estado, sin leer la Tabla 1 de la fuente.

\begin{hafocommonfigure}[H]
\centering
\begin{subfigure}[t]{0.49\textwidth}
  \centering
  \includegraphics[width=\linewidth]{mavpd_nyquist_df.png}
  \caption{Ramas de Nyquist y función descriptiva derivadas.}
\end{subfigure}\hfill
\begin{subfigure}[t]{0.49\textwidth}
  \centering
  \includegraphics[width=\linewidth]{mavpd_python_clusters.png}
  \caption{Dos ciclos internos y candidato externo en Python.}
\end{subfigure}
\caption{Ruta directa MAVPD con \(\xi=3.1\).}
\label{int:fig:mavpd-route}
\end{hafocommonfigure}

\subsection{Comparación cruzada del caso periódico}

\ifHAFOBasinMaterial
El conjunto común contiene \MAVPDProbeCount{} condiciones: 3 equilibrios, 3 radios
\((10^{-5},10^{-3},10^{-2})\) y 12 direcciones explícitas por equilibrio y
radio. Se usaron \(h=0.005\), transitorio 180 y cola 80.
\fi

El comparador posterior confirmó:

\begin{itemize}
  \item diferencia máxima de frecuencias \(2.31\times10^{-14}\), de
        ganancias \(3.50\times10^{-15}\), de amplitudes
        \(3.44\times10^{-15}\) y de semillas \(7.45\times10^{-15}\);
\ifHAFOBasinMaterial
  \item \MAVPDFullAssignmentMatches{} de \MAVPDProbeCount{} asignaciones de
        clúster, clases finales y decisiones sobre el candidato externo
        coincidentes;
  \item 54 trayectorias hacia el ciclo interno positivo, 54 hacia el negativo
        y cero hacia el externo en cada implementación;
\fi
  \item distancia normalizada entre nubes externas
        \(\MAVPDFullCloudDistance\);
  \item periodos \(\MAVPDPythonPeriod\) en Python y
        \(\MAVPDJuliaFullPeriod\) en Julia;
  \item ninguna ruta alternativa para \(\xi=3.1\).
\end{itemize}

\begin{hafocommonfigure}[H]
\centering
\includegraphics[width=0.94\textwidth]{mavpd_python_julia_overlay_full.png}
\caption{Comparación posterior MAVPD \emph{full}. Las localizaciones se
ejecutaron independientemente.
\ifHAFOBasinMaterial
Las 108 coordenadas se compartieron sólo para el control de ocultedad.
\fi}
\label{int:fig:mavpd-overlay}
\end{hafocommonfigure}

\subsection{Clasificación dinámica y discrepancia con la fuente}

La sección de Poincaré Python se refinó con DOP853 y eventos densos:
405 cruces, rango de los últimos puntos del orden de \(10^{-13}\), periodo
\(0.493177986749\) y desviación estándar \(1.25\times10^{-14}\).
El espectro variacional independiente en el cálculo declarado fue
\[
 (-0.0017614,\ -0.3887396,\ -51.1886634),
\]
y la prueba 0--1 dio mediana \(K=0.000440085726\) sobre 64 valores de \(c\).
En conjunto con la sección de Poincaré, estos diagnósticos apoyan un ciclo
estable de periodo uno bajo el protocolo numérico; la prueba 0--1 aislada no
constituye una demostración analítica de periodicidad. El cálculo independiente
en Julia obtuvo \((-0.0042733,-0.3910810,-51.1961919)\), consistente con la
misma clase.

La fuente denomina al conjunto \emph{atractor oculto cuasiperiódico}. Esa
etiqueta no se reproduce para las ecuaciones declaradas: la conclusión local
obtenida es \emph{candidato periódico estable}. La diferencia documenta una
discrepancia reproducible entre la etiqueta de la fuente
y la dinámica obtenida con integración y sección de mayor precisión.

\begin{hafobasinfigure}[H]
\centering
\begin{subfigure}[t]{0.49\textwidth}
  \centering
  \includegraphics[width=\linewidth]{mavpd_python_hiddenness.png}
  \caption{Python: las vecindades alcanzan sólo los ciclos internos.}
\end{subfigure}\hfill
\begin{subfigure}[t]{0.49\textwidth}
  \centering
  \includegraphics[width=\linewidth]{mavpd_julia_candidates.png}
  \caption{Julia \emph{full}: clúster externo con cero impactos.}
\end{subfigure}
\caption{Ciclo externo compatible con ocultedad bajo las 108 sondas del
protocolo finito para \(\xi=3.1\).}
\label{int:fig:mavpd-hiddenness}
\end{hafobasinfigure}

\subsection{Comprobación algebraica de la fuente}

Se preservan tres discrepancias:

\begin{enumerate}
  \item La fuente imprime \(\Theta=-\delta e_1\),
        \(s=-e_1\) y \(\zeta(\sigma)=\sigma^3\). Juntas producen
        \(+\delta y_1^3\), mientras que la ecuación del sistema contiene
        \(-\delta y_1^3\). La formulación evaluada usa la ecuación original y la
        forma exacta \(b=-\delta e_1\), \(c=e_1\).
  \item Para \(\xi=3.1\), los valores impresos
        \(\omega=18.6143\), \(|k|=0.8951\) dejan residuo de cierre
        \(2.45589\) en la transferencia exacta.
  \item Para \(\xi=3.5\), los valores impresos
        \(\omega=10.2344\), \(|k|=0.1136\) dejan residuo \(0.171719\).
\end{enumerate}

Las semillas impresas sí terminan, en una prueba posterior, en los mismos
clústeres dinámicos que la búsqueda independiente. Esa coincidencia dinámica
no elimina la discrepancia algebraica de frecuencia y ganancia.

\ifHAFOBasinMaterial
\subsection{\texorpdfstring{\(\xi=3.5\): control negativo de ocultedad}{xi=3.5: control negativo de ocultedad}}

La ruta directa recalculó dos ramas:
\[
(\omega_0,k,a_0)=
(10.8359213547,\ 0.1497634046,\ 0.4468607607)
\]
y
\[
(\omega_0,k,a_0)=
(13.0511611895,\ 0.3009508811,\ 0.6334570558).
\]
Las cuatro fases convergen a dos ciclos internos simétricos. Bajo el mismo
protocolo de 108 sondeos, Python y Julia observaron 24 contactos
con cada ciclo; el comparador confirmó la misma decisión negativa. Después
del primer contacto se habilitó la ruta alternativa, que no encontró clústeres
nuevos. Por tanto, la etiqueta oculta no se reproduce para \(\xi=3.5\) bajo
las ecuaciones y controles ensayados.

\fi

\subsection{Extensión caótica: continuación relativa a una frontera candidata de Hopf}
\label{int:subsec:mavpd-chaos}

\ifHAFOBasinMaterial
Los dos regímenes anteriores evalúan conjuntos periódicos bajo un protocolo
común.
\else
El régimen periódico anterior establece el punto de comparación dinámico.
\fi
La coexistencia de caos y ocultedad se estudió mediante dos cálculos
independientes, en Python y Julia. Ambos parten de \eqref{int:eq:mavpd},
obtienen las dos semillas directas y recorren 21 niveles de continuación.
Los parámetros y las semillas publicados se reservan para el contraste
posterior.

En el punto base $(\xi,\gamma)=(3.1,0.1)$, los espectros variacionales cortos
de las ramas 0 y 1 fueron, respectivamente,
\[
 (-0.0012506,-1.0123389,-13.1096006),\qquad
 (-0.0029948,-0.3930260,-51.1769818).
\]
Por tanto, ninguna rama superó el cribado finito declarado
$\widehat\lambda_1>0.02$. Este resultado habilita la ruta alternativa; el cribado de tiempo finito
no establece regularidad asintótica.

La ruta alternativa selecciona, entre las ramas integradas, la de
menor frecuencia armónica directa. La regla eligió aquí la rama 0,
$\omega_0=10.5975230$; su estado final de continuación $\lambda$ alimentó el
transporte en parámetros. La rama 1 se evalúa en el punto base y queda pendiente su transporte por A1.

En $E_\pm=(\pm\sqrt{\gamma},0,\pm\sqrt{\gamma})$, el Jacobiano de las
ecuaciones activas produce
\begin{equation}
 p(\lambda)=\lambda^3+a_1\lambda^2+a_2\lambda+a_3,
 \label{int:eq:mavpd-characteristic}
\end{equation}
con
\begin{equation}
 a_1=\xi+2\delta\gamma,\qquad
 a_2=\rho-\delta+2\delta\gamma\xi,\qquad
 a_3=2\delta\gamma\rho.
 \label{int:eq:mavpd-routh-coefficients}
\end{equation}
Al imponer la igualdad de cruce imaginario de Routh--Hurwitz
$a_1a_2=a_3$ y escribir \(g=2\delta\gamma\), la diferencia se desarrolla como
\[
 a_1a_2-a_3=(\xi+g)(\rho-\delta+\xi g)-\rho g
 =\xi g^2+(\xi^2-\delta)g+\xi(\rho-\delta).
\]
Por tanto, la frontera satisface
\begin{equation}
 \xi g^2+(\xi^2-\delta)g+\xi(\rho-\delta)=0.
 \label{int:eq:mavpd-hopf}
\end{equation}
Esta igualdad determina una frontera candidata de Hopf; aquí no se afirma
haber verificado transversalidad ni no degeneración. Para
$\xi=\MAVPDChaosXi$, la raíz alta de \eqref{int:eq:mavpd-hopf} es
$\gamma_H=\MAVPDChaosHopf$. El valor $\xi=\MAVPDChaosXi$ es el extremo local
declarado de la ruta, no una coordenada elegida por el cribado. El procedimiento
transporta primero el estado desde $\xi=3.10$ hasta ese extremo en pasos de
$-0.01$ y después continúa
$\gamma$, reconstruyendo las ecuaciones, el Jacobiano y la forma Lur'e en cada
nodo. La inspección se restringe a desplazamientos declarados respecto de
$\gamma_H$; no es un barrido de frecuencias ni una exploración arbitraria de
condiciones iniciales.

\ifHAFOBasinMaterial
\begin{table}[H]
\centering
\caption{Cribado local MAVPD alrededor de la frontera alta candidata de Hopf. Los
contactos se cuentan sobre 12 sondeos preliminares alrededor de $E_0$; la tabla
redondea $\widehat\lambda_1$ a un máximo de seis decimales, pero los cálculos
conservan la precisión completa.}
\label{int:tab:mavpd-chaos-screen}
\small
\begin{tabular}{@{}C{1.8cm}C{2.0cm}C{2.3cm}L{8.7cm}@{}}
\toprule
$\gamma-\gamma_H$ & $\widehat\lambda_1$ & Contactos desde $E_0$ & Lectura del cribado \\
\midrule
% Generado desde todas las filas de 03_candidate_screening.csv.
% Columnas: hopf_offset; lambda_1; E0_target_hits/E0_probe_count; lectura finita.
0.002 & 0.929112 & 6/12 & No elegible: hubo contactos en los sondeos preliminares. \\
0.003 & 0.893403 & 6/12 & No elegible: hubo contactos en los sondeos preliminares. \\
0.005 & 0.188317 & 6/12 & No elegible: hubo contactos en los sondeos preliminares. \\
0.008 & 0.812209 & 0/12 & Elegible bajo el cribado finito; no seleccionada por la regla. \\
0.010 & 0.72394 & 0/12 & Elegible y seleccionada por la regla declarada. \\
0.012 & 0.315116 & 0/12 & No elegible bajo el cribado finito declarado. \\
0.015 & -0.00887 & 0/12 & No elegible bajo el cribado finito declarado. \\
\bottomrule

\end{tabular}
\end{table}
\fi

La regla selecciona el mayor desplazamiento declarado con cálculo de Lyapunov
válido y $\widehat\lambda_1>0.5$
\ifHAFOBasinMaterial
, cero contactos preliminares, cero casos ambiguos y cero fallos numéricos
\else
y con los demás filtros de admisibilidad satisfechos
\fi
: así se obtiene
\[
 (\xi,\gamma,\delta,\rho)=
 (\MAVPDChaosXi,\MAVPDChaosGamma,100,200),\qquad
 \gamma=\gamma_H+0.010.
\]
El desplazamiento $0.010$ es el resultado de la regla de selección aplicada en
esta continuación numérica local; no se lee como valor esperado ni es un
parámetro publicado en \cite{matouk2023}. Por tanto, sólo $\gamma$ se selecciona
en este cribado: la tupla completa no se presenta como una optimización conjunta
de $\xi$ y $\gamma$.

\begin{hafobasinfigure}[H]
\centering
\includegraphics[width=0.90\textwidth]{mavpd_chaos_screen.png}
\caption{Continuación en parámetros y cribado conjunto de exponente máximo y
contactos desde $E_0$. La selección ocurre después del fallo dinámico de las
ramas directas, no mediante una malla de frecuencias.}
\label{int:fig:mavpd-chaos-screen}
\end{hafobasinfigure}

\begin{table}[H]
\centering
\caption{Protocolo numérico de la extensión MAVPD caótica.}
\label{int:tab:mavpd-chaos-contract}
\footnotesize
\begin{tabularx}{\textwidth}{@{}L{0.19\textwidth}L{0.34\textwidth}Y@{}}
\toprule
Fase & Horizonte y muestreo & Integrador y tolerancias \\
\midrule
Cribado D0 por rama & $t_b=100$, $t_{LE}=250$, QR $0.5$ &
DOP853, $(\mathrm{rtol},\mathrm{atol})=(2\,10^{-10},2\,10^{-12})$,
$\Delta t_{\max}=0.02$ \\
Transporte A1 por nodo & $t_b=60$, $t_k=10$, muestra $0.02$ &
DOP853, $(2\,10^{-10},2\,10^{-12})$, $\Delta t_{\max}=0.02$ \\
Referencia del cribado & duración 360, descarte 160, muestra $0.05$,
600 puntos máximos & DOP853, $\Delta t_{\max}=0.03$; factor de seguridad 3 \\
\ifHAFOBasinMaterial
Sondeos preliminares $E_0$ & $t_b=500$, $t_k=100$, muestra $0.05$;
2 radios, 6 direcciones & DOP853, $(2\,10^{-10},2\,10^{-12})$,
$\Delta t_{\max}=0.03$ \\
\fi
Trayectoria estricta & $t_f=900$, $t_b=300$, muestra $0.02$ &
DOP853, $(2\,10^{-11},2\,10^{-13})$, $\Delta t_{\max}=0.01$ \\
Lyapunov estricto & $t_b=300$, $t_{LE}=1200$, QR $0.5$ &
DOP853 variacional, $(2\,10^{-12},2\,10^{-14})$,
$\Delta t_{\max}=0.01$ \\
Control Lyapunov combinado & estado candidato perturbado; $t_b=200$,
$t_{LE}=600$, QR $0.5$ & DOP853 variacional,
$(2\,10^{-10},2\,10^{-12})$, $\Delta t_{\max}=0.02$ \\
\ifHAFOBasinMaterial
Ocultedad final & $t_b=400$, $t_k=100$, muestra $0.05$;
108 sondeos principales + 4 dirigidos & DOP853, $(2\,10^{-10},2\,10^{-12})$,
$\Delta t_{\max}=0.03$; referencia máxima 1000 puntos \\
\fi
\bottomrule
\end{tabularx}
\end{table}

\subsubsection*{Refinamiento dinámico y controles de caos}

La trayectoria candidata se integró con DOP853, $t_f=900$, transitorio 300,
muestreo $0.02$, tolerancias $(\mathrm{rtol},\mathrm{atol})=(2\times10^{-11},
2\times10^{-13})$ y paso máximo $0.01$ \cite{scipy2020}. Sobre un problema
variacional separado, DOP853 con QR, transitorio 300 y acumulación 1200 produjo
el espectro finito. Con $Q_0=I$, en cada segmento se integra
\begin{equation}
 \dot Y_m=J(y(t))Y_m,\qquad Y_m(t_{m-1})=Q_{m-1},
 \qquad Y_m(t_m)=Q_mR_m,
 \quad
 \widehat\lambda_i(T)=\frac{1}{T}\sum_{m=1}^{M}
 \log\left|R_{m,ii}\right|,
 \label{int:eq:mavpd-variational-qr}
\end{equation}
donde cada $R_m$ contiene sólo el crecimiento del segmento correspondiente;
la base se reinicia con $Q_m$ antes del segmento siguiente
\cite{Benettin1980a}. El resultado fue
\begin{equation}
 \begin{aligned}
  \widehat\lambda_1&=\MAVPDChaosLEOne,
  &\widehat\lambda_2&=\MAVPDChaosLETwo,\\
  \widehat\lambda_3&=\MAVPDChaosLEThree,
  &D_{KY}&=\MAVPDChaosKY.
 \end{aligned}
 \label{int:eq:mavpd-chaos-lyapunov}
\end{equation}
Aquí \(D_{KY}=2+(\widehat\lambda_1+\widehat\lambda_2)/
|\widehat\lambda_3|\), pues las dos primeras suman positivamente y la suma de
las tres es negativa.
El control de Lyapunov perturbó el extremo de la continuación y cambió simultáneamente las
tolerancias, el paso máximo y los horizontes de transitorio y acumulación. El cálculo conservó $\widehat\lambda_1=\MAVPDChaosLEControlOne$. El residuo de
consistencia
$r_{\mathrm{div}}=\sum_i\widehat\lambda_i(T_{LE})-
\overline{\nabla\!\cdot f}_{[t_a,t_b]}$ fue
$\MAVPDChaosDivergenceResidual$, con
\begin{equation}
 \nabla\!\cdot f(y)=\delta\gamma-\xi-3\delta y_1^2,
 \qquad
 \overline{\nabla\!\cdot f}_{[t_a,t_b]}=\frac{1}{t_b-t_a}
 \int_{t_a}^{t_b}\!\left(\delta\gamma-\xi-3\delta y_1(t)^2\right)\,dt.
 \label{int:eq:mavpd-divergence-control}
\end{equation}
La relación con el espectro se obtiene de la fórmula de Jacobi. Para una
matriz fundamental invertible \(\Phi\),
\[
 \frac{\dd}{\dd t}\log|\det\Phi(t)|
 =\operatorname{tr}(\Phi^{-1}J\Phi)=\operatorname{tr}J
 =\nabla\!\cdot f.
\]
En la descomposición por segmentos, los determinantes de las matrices
ortogonales \(Q_m\) tienen módulo uno y cada matriz triangular satisface
\(\det R_m=\prod_iR_{m,ii}\). Por tanto, sobre una misma trayectoria y
ventana de longitud \(T\),
\[
 \sum_i\widehat\lambda_i(T)
 =\frac1T\log|\det\Phi(T)|
 =\frac1T\int_0^T\nabla\!\cdot f(y(t))\,\dd t,
\]
con \(\Phi(0)=I\). Esta identidad justifica la comparación de la suma de
exponentes con el promedio de la divergencia.
Aquí $T_{LE}=1200$ pertenece a la integración variacional, mientras que la
divergencia se promedia sobre la trayectoria muestreada independiente
$[t_a,t_b]=[300,900]$. El residuo compara, por tanto, promedios de trayectorias y ventanas
distintas. Como control de integrador, una
trayectoria EFORK de paso fijo $h=0.002$ alcanzó la misma nube de referencia
calibrada.

\begin{hafocommonfigure}[H]
\centering
\begin{subfigure}[t]{0.98\textwidth}
  \centering
  \includegraphics[width=\linewidth]{mavpd_chaos_phase_portraits.png}
  \caption{Proyecciones postransitorias del candidato seleccionado.}
\end{subfigure}

\vspace{0.35em}
\begin{subfigure}[t]{0.86\textwidth}
  \centering
  \includegraphics[width=\linewidth]{mavpd_chaos_time_series.png}
  \caption{Series temporales en la ventana conservada.}
\end{subfigure}
\caption{Geometría y evolución temporal del candidato MAVPD caótico. Ventana postransitoria acotada.}
\label{int:fig:mavpd-chaos-trajectory}
\end{hafocommonfigure}

La sección $y_2=0$ con orientación positiva retuvo
\MAVPDChaosPoincareCount{} cruces no duplicados. Los cruces se obtuvieron por
interpolación lineal entre muestras de la trayectoria, con el error temporal
asociado a esa aproximación
\cite{henon1982}. La prueba 0--1 sobre dos coordenadas de la secuencia de
retorno dio medianas $K=\MAVPDChaosReturnKOne$ y
$K=\MAVPDChaosReturnKTwo$ \cite{gottwald2009}. En cambio, la prueba 0--1 sobre
el flujo fue sensible al submuestreo: 3 de 10 pasos dieron etiqueta de
candidato caótico, 1 dio etiqueta regular y 6 quedaron inconclusos. El criterio conjunto usa la secuencia de retorno con
el paso fijado antes de la clasificación.

\begin{hafocommonfigure}[H]
\centering
\begin{subfigure}[t]{0.49\textwidth}
  \centering
  \includegraphics[width=\linewidth]{mavpd_chaos_lyapunov.png}
  \caption{Convergencia del espectro variacional finito.}
\end{subfigure}\hfill
\begin{subfigure}[t]{0.49\textwidth}
  \centering
  \includegraphics[width=\linewidth]{mavpd_chaos_poincare.png}
  \caption{Sección muestreada e interpolada linealmente.}
\end{subfigure}
\caption{Diagnósticos dinámicos principales del candidato caótico.}
\label{int:fig:mavpd-chaos-diagnostics}
\end{hafocommonfigure}

El espectro de potencia FFT normalizado tiene líneas dominantes y resulta
inconcluso para clasificar el régimen. Se interpreta junto con los
exponentes variacionales y los retornos de Poincaré.

\begin{hafocommonfigure}[H]
\centering
\includegraphics[width=0.82\textwidth]{mavpd_chaos_fft.png}
\caption{Potencia FFT normalizada del candidato.}
\label{int:fig:mavpd-chaos-fft}
\end{hafocommonfigure}

\ifHAFOBasinMaterial
\subsubsection*{Ocultedad bajo vecindades ensayadas}

El clasificador de la nube candidata se calibró con tres ventanas temporales
disjuntas de la misma trayectoria y verificó su separación usando los tres
equilibrios como controles negativos. Para dos nubes finitas \(A,B\), usa la
mediana simétrica de vecinos más cercanos, normalizada por la extensión \(S\)
de la referencia:
\begin{equation}
 d_N(A,B)=\frac{1}{S}\operatorname{mediana}\!\left(
 \{\min_{b\in B}\|a-b\|:a\in A\}\cup
 \{\min_{a\in A}\|b-a\|:b\in B\}\right).
 \label{int:eq:mavpd-cloud-distance}
\end{equation}
Si $R_1,R_2,R_3$ son las ventanas y \(R=R_1\cup R_2\cup R_3\), la escala es
\begin{equation}
 S=\left[\sum_{k=1}^{3}
 \left(\max_{x\in R}x_k-\min_{x\in R}x_k\right)^2\right]^{1/2},
 \qquad
 D(A)=\operatorname{mediana}_{j=1,2,3}d_N(A,R_j).
 \label{int:eq:mavpd-cloud-classifier}
\end{equation}
Si $b$ es el cuantil 0.95 de las distancias entre ventanas, el umbral es
$\tau=3b$ y el margen ambiguo es $m=\max(0.25\tau,b)$. En esta configuración,
$m=b=\tau/3$. Los controles negativos no fijan $\tau$: exigen
$\min_j d_N(N_\ell,R_j)>\tau+m$; de lo contrario se concluye que los controles
se solapan y la calibración se rechaza.

El protocolo principal contiene \MAVPDChaosMainProbes{} trayectorias:
3 equilibrios, radios $(10^{-5},10^{-3},10^{-2})$ y 12 direcciones esféricas
deterministas por radio. Se añadieron \MAVPDChaosTargetedProbes{} controles
dirigidos sobre ambos signos del autovector inestable de $E_0$, en radios
$10^{-7}$ y $10^{-4}$. En total hubo \MAVPDChaosTotalProbes{} trayectorias,
\MAVPDChaosHits{} contactos con la nube candidata, cero ambigüedades y cero
fallos numéricos; todas terminaron en un equilibrio.

\begin{hafobasinfigure}[H]
\centering
\includegraphics[width=0.88\textwidth]{mavpd_chaos_hiddenness.png}
\caption{Resultados de los sondeos alrededor de $E_0$, $E_+$ y $E_-$, junto
con los controles dirigidos del autovector inestable. Cero contactos es una
conclusión finita restringida a este protocolo.}
\label{int:fig:mavpd-chaos-hiddenness}
\end{hafobasinfigure}

La evaluación conjunta clasificó el resultado como compatible con caos y
ocultedad únicamente bajo las vecindades ensayadas. Los diagnósticos numéricos
y la ausencia de contactos se restringen a las ventanas y vecindades
declaradas; la desconexión global de la cuenca permanece sin demostrar.
\fi

\subsubsection*{Comparación numérica independiente}

El cálculo independiente en Julia partió de las ecuaciones y el Jacobiano, recalculó las
dos ramas D0 sin malla de frecuencias y sólo activó A1 después de que ambas
fallaron el umbral dinámico. La regla de menor $\omega_0$ volvió a seleccionar
la rama 0; después transportó el estado de $\xi=3.10$ a $2.85$, derivó la misma
frontera de Routh--Hurwitz y seleccionó de nuevo el desplazamiento $0.010$.
Durante esa localización no leyó la semilla, el estado terminal ni $\gamma$ de
Python.
\ifHAFOBasinMaterial
El intercambio de las condiciones de ocultedad ocurrió únicamente en la fase
posterior de comparación.
\fi

Python empleó DOP853 adaptativo para la trayectoria y la ecuación
variacional. Julia empleó RK4 de paso fijo para las trayectorias y los sondeos,
y reortogonalización QR para el cálculo variacional.

\begin{table}[H]
\centering
\caption{Comparación del candidato MAVPD caótico con tolerancias comunes.}
\label{int:tab:mavpd-chaos-cross-runtime}
\small
\begin{tabularx}{\textwidth}{@{}L{0.31\textwidth}YY@{}}
\toprule
Magnitud & Python/\hafo{} & Julia \\
\midrule
$\xi$ & \MAVPDChaosXi & \MAVPDChaosXi \\
$\gamma$ & \MAVPDChaosGamma & \MAVPDChaosGamma \\
$\widehat\lambda_1$ & \MAVPDChaosLEOne & \MAVPDChaosJuliaLEOne \\
$\widehat\lambda_2$ & \MAVPDChaosLETwo & \MAVPDChaosJuliaLETwo \\
$\widehat\lambda_3$ & \MAVPDChaosLEThree & \MAVPDChaosJuliaLEThree \\
$D_{KY}$ & \MAVPDChaosKY & \MAVPDChaosJuliaKY \\
Cruces de Poincaré & \MAVPDChaosPoincareCount & \MAVPDChaosJuliaPoincareCount \\
Mediana $K$ sobre retornos & \MAVPDChaosReturnKOne & \MAVPDChaosJuliaReturnK \\
\ifHAFOBasinMaterial
Decisiones de sondeo coincidentes &
\MAVPDChaosTotalProbes/\MAVPDChaosTotalProbes &
\MAVPDChaosJuliaProbeMatches/\MAVPDChaosTotalProbes \\
\fi
\bottomrule
\end{tabularx}
\end{table}

La comparación cruzada coincidió exactamente en $\xi$, $\gamma$ y el
desplazamiento de Hopf. Las estadísticas de la nube postransitoria y los
diagnósticos de caos quedaron dentro de sus bandas predeclaradas.
\ifHAFOBasinMaterial
Además, coincidieron las
\MAVPDChaosTotalProbes{} condiciones iniciales, estados, destinos y decisiones
de contacto.
\fi

\begin{hafobasinfigure}[H]
\centering
\begin{subfigure}[t]{0.49\textwidth}
  \centering
  \includegraphics[width=\linewidth]{julia_mavpd_chaos_gamma_screen.png}
  \caption{Cribado y selección de $\gamma$ recalculados en Julia.}
\end{subfigure}\hfill
\begin{subfigure}[t]{0.49\textwidth}
  \centering
  \includegraphics[width=\linewidth]{julia_mavpd_chaos_hiddenness_outcomes.png}
  \caption{Decisión de los 112 sondeos finitos en Julia.}
\end{subfigure}
\caption{Artefactos de la reproducción Julia MAVPD. El primer panel es un
diagnóstico de continuación; el segundo resume un protocolo finito de cuenca.}
\label{int:fig:mavpd-chaos-julia}
\end{hafobasinfigure}

\begin{hafocommonfigure}[H]
\centering
\begin{subfigure}[t]{0.49\textwidth}
  \centering
  \includegraphics[width=\linewidth]{julia_mavpd_chaos_nyquist_diagnostic.png}
  \caption{Nyquist de la transferencia ya derivada.}
\end{subfigure}\hfill
\begin{subfigure}[t]{0.49\textwidth}
  \centering
  \includegraphics[width=\linewidth]{julia_mavpd_chaos_amplitude_spectrum.png}
  \caption{Espectro de amplitud FFT postransitorio.}
\end{subfigure}
\caption{Diagnósticos intermedios Julia MAVPD. Nyquist no suministra una malla
de búsqueda ni selecciona $\omega_0$; la FFT no se denomina PSD.}
\label{int:fig:mavpd-chaos-julia-spectral}
\end{hafocommonfigure}


\section{Caso 4: PLL bifásico con filtro lead--lag}

\subsection{Transformación T1 y obstrucción directa}

El modelo publicado es
\begin{align}
 \dot x&=-\frac{x}{T}+\frac{\tau_1}{2T}\sin\theta,\nonumber\\
 \dot\theta&=\omega_\Delta-\frac{Lx}{T}
 -\frac{\tau_2L}{2T}\sin\theta,\qquad T=\tau_1+\tau_2,
 \label{int:eq:pll}
\end{align}
con \(\tau_1=0.0448\), \(\tau_2=0.0185\), \(L=500\) y
\(\omega_\Delta=178.9\) \cite{pll2015}. El desplazamiento de un equilibrio
bloqueado elimina el término afín, de modo que el sistema es T1. Como
\(\theta\) es angular, la recurrencia y las distancias se calculan módulo
\(2\pi\), no con una norma euclidiana sin envolver.

En un equilibrio, la primera ecuación da
\(x_e=(\tau_1/2)\sin\theta_e\). Sustituir en la segunda y usar
\(T=\tau_1+\tau_2\) conduce a
\[
 0=\omega_\Delta-\frac{L(\tau_1+\tau_2)}{2T}\sin\theta_e,
 \qquad \sin\theta_e=\frac{2\omega_\Delta}{L}=0.7156.
\]
En una vuelta del cilindro hay dos soluciones:
\[
 x_e=0.01602944,\qquad
 \theta_f=0.797482699882,\qquad
 \theta_s=2.344109953708.
\]
El Jacobiano en cualquiera de ellas es
\[
 J_e=\begin{pmatrix}
 -1/T&\tau_1\cos\theta_e/(2T)\\
 -L/T&-\tau_2L\cos\theta_e/(2T)
 \end{pmatrix},
\]
con
\[
 \operatorname{tr}J_e=-\frac{1+\tau_2L\cos\theta_e/2}{T},\qquad
 \det J_e=\frac{L\cos\theta_e}{2T}.
\]
En \(\theta_f\), la traza es negativa, el determinante positivo y el
discriminante \((\operatorname{tr}J_e)^2-4\det J_e\) es negativo para estos
parámetros: se obtiene un foco estable. En \(\theta_s\), el determinante es
negativo y el equilibrio es una silla.

El cambio \(u=x-x_e\), \(v=\theta-\theta_e\) transforma el campo en
\[
 \begin{pmatrix}\dot u\\\dot v\end{pmatrix}
 =P\begin{pmatrix}u\\v\end{pmatrix}+b\psi(v),
\]
\[
 P=\begin{pmatrix}-1/T&0\\-L/T&0\end{pmatrix},\quad
 b=\begin{pmatrix}\tau_1/(2T)\\-\tau_2L/(2T)\end{pmatrix},
 \quad c=\begin{pmatrix}0\\1\end{pmatrix},
\]
con \(\psi(v)=\sin(v+\theta_e)-\sin\theta_e\). Las identidades del
equilibrio cancelan los términos constantes. Para una entrada \(\psi\),
la primera fila del resolvente \((sI-P)^{-1}b\) da
\[
 u=\frac{\tau_1}{2(1+Ts)}\psi,
\]
y la segunda da
\[
 sv=-\frac L{2T}\left[\frac{\tau_1}{1+Ts}+\tau_2\right]\psi
    =-\frac L2\frac{1+\tau_2s}{1+Ts}\psi.
\]
Por tanto, las dos convenciones de transferencia se relacionan mediante
\begin{equation}
 W(s)=c^{\mathsf T}(sI-P)^{-1}b
     =-\frac L2\frac{1+\tau_2s}{s(1+Ts)}=-G(s).
 \label{int:eq:pll-transfer}
\end{equation}
Al multiplicar el cociente en \(s=i\omega\) por el conjugado de su
denominador,
\[
 W(i\omega)=\frac{L\tau_1}{2(1+T^2\omega^2)}
 +i\,\frac{L(1+\tau_2T\omega^2)}{2\omega(1+T^2\omega^2)}.
\]
Así, \(\operatorname{Im}W(i\omega)>0\) y
\(\operatorname{Im}G(i\omega)<0\) para todo \(\omega>0\). Ninguna
convención admite el cruce real requerido por una ganancia descriptiva real;
la ruta directa no produce candidatos.

\subsection{Ruta A1 de Andronov}

Para \(L>0\), definir \(y=\dot\theta\) permite recuperar la variable
original mediante
\begin{equation}
 x=\frac T L(\omega_\Delta-y)-\frac{\tau_2}{2}\sin\theta.
 \label{int:eq:pll-inverse-andronov}
\end{equation}
Derivando la segunda ecuación de \eqref{int:eq:pll},
\[
 \dot y=-\frac L T\dot x-\frac{\tau_2L}{2T}\cos\theta\,y
        =\frac L{T^2}x-\frac{L\tau_1}{2T^2}\sin\theta
         -\frac{\tau_2L}{2T}\cos\theta\,y.
\]
Sustituir \eqref{int:eq:pll-inverse-andronov} y agrupar
\(\tau_1+\tau_2=T\) produce
\begin{equation}
 T\dot y=\omega_\Delta-\frac{L}{2}\sin\theta
 -\left(1+\frac{\tau_2L}{2}\cos\theta\right)y.
 \label{int:eq:pll-andronov}
\end{equation}
Para una rotación positiva se exige \(y(\theta)>0\) durante toda la
vuelta. Entonces \(\theta\) puede usarse como variable independiente:
\begin{equation}
 \frac{\dd y}{\dd\theta}=H_L(\theta,y)
 =\frac{\omega_\Delta-(L/2)\sin\theta}{Ty}
  -\frac{1+(\tau_2L/2)\cos\theta}{T}.
 \label{int:eq:pll-angle-flow}
\end{equation}
Con \(y(0)=y_0>0\), el mapa de retorno es
\(P_L(y_0)=y(2\pi;y_0)\). Una órbita rotatoria cumple
\(P_L(y_0)-y_0=0\), y su periodo físico es
\[
 \mathcal T(y_0)=\int_0^{2\pi}\frac{\dd\theta}{y(\theta;y_0)}.
\]
La derivada del retorno se calcula diferenciando la ecuación respecto de
\(y_0\). Si \(z=\partial y/\partial y_0\),
\[
 \frac{\dd z}{\dd\theta}
 =-\frac{\omega_\Delta-(L/2)\sin\theta}{Ty^2}z,
 \qquad z(0)=1,\qquad P_L'(y_0)=z(2\pi).
\]
El multiplicador transversal es \(P_L'(y_0)\): módulo menor que uno indica
estabilidad local, y módulo mayor que uno inestabilidad. El mapa queda
indefinido para esta construcción si la solución pierde la rotación positiva
antes de completar \(2\pi\).

En \(L=0\), la ecuación reducida admite \(y(\theta)\equiv\omega_\Delta\),
con
\[
 \mathcal T_0=\frac{2\pi}{\omega_\Delta},\qquad
 P_0'(\omega_\Delta)=\exp\!\left(-\frac{2\pi}{T\omega_\Delta}\right).
\]
En las coordenadas originales, \(\theta(t)=\omega_\Delta t\) y el filtro
satisface una ecuación lineal forzada. Una solución particular periódica
\(x_p=A\sin\theta+B\cos\theta\) cumple
\(A-\omega_\Delta TB=\tau_1/2\) y
\(\omega_\Delta TA+B=0\). Por tanto,
\[
 x_p(\theta)=\frac{\tau_1}{2[1+(\omega_\Delta T)^2]}
             [\sin\theta-\omega_\Delta T\cos\theta].
\]
Las restantes soluciones difieren de ella en \(C\e^{-t/T}\), lo cual
concuerda con el multiplicador transversal calculado.
Esta solución inicia la continuación de la ecuación reducida en 101 valores
hasta \(L=500\). La reconstrucción de \(x\) mediante
\eqref{int:eq:pll-inverse-andronov} se aplica sólo en los niveles \(L>0\),
pues el cambio de variables es singular en \(L=0\). Las condiciones iniciales
publicadas se reservan para la comparación posterior.

El análisis localizó:

\begin{table}[H]
\centering
\caption{Órbitas de retorno PLL en \(L=500\).}
\small
\begin{tabular}{@{}lrrrr@{}}
\toprule
Órbita & \(y_0\) & \(x_0\) & Periodo & Multiplicador \\
\midrule
Ciclo estable &
140.125084913 & 0.004908904250 &
\PLLPythonStablePeriod & \PLLPythonStableMultiplier \\
Separatriz inestable &
135.172047420 & 0.005535958797 &
\PLLPythonSeparatrixPeriod & \PLLPythonSeparatrixMultiplier \\
\bottomrule
\end{tabular}
\end{table}

Los errores de retorno Python sobre el cilindro fueron
\(2.68\times10^{-15}\) para el ciclo estable y
\(6.18\times10^{-14}\) para la separatriz. Julia reprodujo los periodos
\(\PLLJuliaStablePeriod\) y \(\PLLJuliaSeparatrixPeriod\), con multiplicadores
\(\PLLJuliaStableMultiplier\) y \(\PLLJuliaSeparatrixMultiplier\). Las tres
condiciones iniciales publicadas quedaron fuera del flujo de localización: el
control posterior confirmó que dos alcanzan el ciclo y una el foco.

\begin{hafocommonfigure}[H]
\centering
\begin{subfigure}[t]{0.49\textwidth}
  \centering
  \includegraphics[width=\linewidth]{pll_transfer_obstruction.png}
  \caption{Obstrucción analítica de la ruta directa.}
\end{subfigure}\hfill
\begin{subfigure}[t]{0.49\textwidth}
  \centering
  \includegraphics[width=\linewidth]{pll_phase_cylinder_cycles.png}
  \caption{Ciclo estable y separatriz en \(\mathbb R\times S^1\).}
\end{subfigure}
\caption{Localización T1 + A1 del PLL lead--lag.}
\label{int:fig:pll-route}
\end{hafocommonfigure}

\ifHAFOBasinMaterial
\subsection{Sondeos cilíndricos y estado comparativo}

El sondeo común usó \PLLProbeCount{} condiciones:
2 equilibrios, 4 radios \((10^{-5},10^{-4},10^{-3},10^{-2})\) y 12
direcciones por radio. Ninguna alcanzó el ciclo. Las 48 condiciones alrededor
del foco y las 48 alrededor del punto silla terminaron en el foco. Por tanto,
el resultado es compatible con ocultedad bajo el sondeo cilíndrico declarado.

\begin{hafobasinfigure}[H]
\centering
\includegraphics[width=0.82\textwidth]{julia_pll_cylindrical_hiddenness_probes.png}
\caption{Visualización Julia de las condiciones iniciales comunes del control
PLL, expresadas en las coordenadas envueltas del cilindro.}
\label{int:fig:pll-hiddenness}
\end{hafobasinfigure}
\fi

Julia implementó independientemente la misma distancia angular, el mapa de
retorno y la continuación desde la solución exacta en $L=0$; no leyó estados
Python ni las condiciones publicadas para localizar las órbitas.
\ifHAFOBasinMaterial
Coincidieron
\PLLJuliaProbeMatches{} de \PLLProbeCount{} estados y decisiones de contacto,
así como los destinos de las \PLLJuliaDestinationMatches{} trayectorias.
\fi
Las órbitas interpoladas sobre una fase normalizada difirieron como máximo
\PLLStableOrbitMaxDistance{} (estable) y
\PLLSeparatorOrbitMaxDistance{} (separatriz) en la métrica cilíndrica escalada.


\begin{hafocommonfigure}[H]
\centering
\begin{subfigure}[t]{0.49\textwidth}
  \centering
  \includegraphics[width=\linewidth]{julia_pll_loop_gain_continuation.png}
  \caption{Continuación Julia desde $L=0$ hasta $L=500$.}
\end{subfigure}\hfill
\begin{subfigure}[t]{0.49\textwidth}
  \centering
  \includegraphics[width=\linewidth]{julia_pll_phase_cylinder_cycles.png}
  \caption{Ciclo estable y separatriz Julia en el cilindro.}
\end{subfigure}
\caption{Continuación y órbitas de la ruta PLL T1 + A1 calculadas en Julia.}
\label{int:fig:pll-julia}
\end{hafocommonfigure}

\section{Sistemas enteros estudiados y pendientes}

\subsection{Taxonomía D0/T1/A1/M1}

\begin{table}[H]
\centering
\caption{Clases estructurales usadas para aceptar o posponer sistemas.}
\small
\begin{tabularx}{\textwidth}{@{}C{1.1cm}p{0.24\textwidth}Yp{0.25\textwidth}@{}}
\toprule
Clase & Definición & Qué permite & Restricción \\
\midrule
D0 & Forma escalar directa \(\dot X=PX+b\psi(c^{\mathsf T}X)\) &
Ruta de transferencia y función descriptiva sin cambio de variables. &
Un solo canal no lineal escalar. \\
T1 & Transformación compatible e invertible, desplazamiento de equilibrio,
variedad invariante documentada o PWA de rango uno &
Lleva el sistema a D0 conservando la interpretación dinámica necesaria. &
La transformación registra la inversa y la métrica correspondiente, por ejemplo angular. \\
A1 & La forma escalar existe, pero la ruta directa se rechaza con una razón
matemática &
Habilita mapa de conmutación, homotopía publicada u otra alternativa. &
Conserva el diagnóstico de incompatibilidad directa y excluye semillas
publicadas de la inicialización. \\
M1 & Varios canales no lineales independientes &
Reserva para una futura extensión multicanal. &
Requiere una formulación multicanal adicional. \\
\bottomrule
\end{tabularx}
\end{table}

\subsection{Cobertura de los casos de estudio}

La \cref{int:tab:catalog} distingue los sistemas integrados de los candidatos
cuya formulación o evaluación numérica está pendiente.

\begin{longtable}{@{}C{1.2cm}L{0.29\textwidth}C{1.65cm}L{0.42\textwidth}@{}}
\caption{Sistemas de prueba de orden entero.}
\label{int:tab:catalog}\\
\toprule
Prior. & Sistema & Clase & Estado y papel científico \\
\midrule
\endfirsthead
\toprule
Prior. & Sistema & Clase & Estado y papel científico \\
\midrule
\endhead
Base & Chua no suave entero & D0 &
Compatible con ocultedad bajo el protocolo finito; reproducido en Python y Julia. \\
Base & Kalman--Fitts de cuarto orden & D0 + A1 &
No Chua; ciclo límite localizado por mapa de conmutación y continuación,
compatible con ocultedad bajo el sondeo; reproducido en Python y Julia. \\
1 & Van der Pol--Duffing autónomo modificado & D0 &
\(\xi=3.1\): 108/108 coincidencias entre cálculos y cero
contactos externos.
\ifHAFOBasinMaterial
\(\xi=3.5\): control negativo, 24/24 contactos con el ciclo seleccionado.
\fi
\(\xi=\MAVPDChaosXi\): candidato caótico compatible con
ocultedad bajo \MAVPDChaosTotalProbes{} sondeos en las vecindades declaradas;
coincidencia Python--Julia en todos los destinos evaluados. \\
2 & PLL bifásico con filtro lead--lag & T1 + A1 &
Comparación Python--Julia con 96 sondeos; las órbitas, su estabilidad
y las decisiones coinciden. \\
3 & Leonov--Kuznetsov, sistema (162) & D0 + A1 &
Ecuaciones y ruta de saturación disponibles; falta verificar la fórmula
completa por tramos y fijar los parámetros de integración. \\
4 & Jerk exponencial de Fischer, \(q=1\) & D0 con balance DC &
Control bibliográfico y de Lyapunov: equilibrio inestable; se conserva la
discrepancia del tercer exponente. Sin una campaña de cuenca local, la evidencia
disponible no permite clasificarlo como autoexcitado ni oculto. \\
\bottomrule
\end{longtable}

\subsection{Siguientes casos y reservas}

Leonov--Kuznetsov, sistema (162), es el siguiente candidato D0 + A1. La
fuente usa una familia de saturaciones, transporta el estado terminal entre
etapas y luego continúa la no linealidad hasta \(\tanh\)
\cite{LeonovKuznetsov2013}. Su evaluación requiere verificar la fórmula
completa por tramos y fijar los parámetros de integración.

El jerk exponencial de Fischer se conserva como control D0 con balance DC y
como banco de exponentes con discrepancias documentadas \cite{fischer2020}.
Falta un sondeo de vecindades para evaluar
la clasificación de cuenca. La inestabilidad del equilibrio, por sí sola, no
implica autoexcitación. La evaluación entera requiere los parámetros del caso bibliográfico correspondiente.

Quedan en reserva Genesio--Tesi, un oscilador memristivo cúbico, el jerk
hiperbólico de Joshi--Ranjan, un jerk PWA y los jerks de Rech y
Rasul--Salih. Faltan parámetros, condiciones iniciales o la verificación de
la ecuación exacta en su fuente primaria.

Lorenz, Rabinovich--Fabrikant, el sistema financiero, el sistema \emph{four
wing} y los modelos con productos no lineales independientes son M1. Su estudio
con Julia queda fuera de las hipótesis de la ruta Lur'e escalar D0/T1.

\section{Comparación metodológica}

\begin{table}[H]
\centering
\caption{Características de los métodos de localización y comparación.}
\small
\begin{tabularx}{\textwidth}{@{}p{0.20\textwidth}YY@{}}
\toprule
Aspecto & Python/\hafo{} & Julia/\juliads--\attractorsjl{} \\
\midrule
Entrada científica & Ecuaciones, Jacobiano, equilibrios, forma Lur'e,
transferencia y función descriptiva & Ecuaciones y dominio para búsqueda
global; las mismas declaraciones Lur'e cuando se replica la ruta teórica \\
Localización & Dirigida por balance armónico y continuación & Global por
recurrencias en Chua; ruta teórica equivalente en Kalman--Fitts y en el MAVPD
periódico y caótico; ruta A1 cilíndrica equivalente en PLL \\
Frecuencia y semilla & Produce \(\omega_0,k,a_0,X_0\) con residuos & Puede
obtener esos valores mediante la ruta teórica; la búsqueda por recurrencias
no utiliza ese balance \\
Generalidad & Restringida a D0/T1 y a las alternativas A1 registradas &
Más amplia para mapas, flujos, cuencas y recurrencias \\
Ocultedad & Flujo explícito de equilibrios y nube candidata & Protocolo
añadido al ejemplo; no es una etiqueta automática de la biblioteca \\
Valores publicados & Regresión posterior & Regresión posterior \\
Fallo típico & No hay ganancia/amplitud compatible o se pierde la rama &
La cuenca no es muestreada, la partición no resuelve el conjunto o el
clasificador agrupa mal \\
\bottomrule
\end{tabularx}
\end{table}

La forma Lur'e relaciona explícitamente el campo, la transferencia y la semilla. La búsqueda por recurrencias aporta generalidad y no necesita una
transformación Lur'e, pero
su éxito depende del dominio, partición, número de muestras y tamaño de la
cuenca. La comparación de los candidatos requiere un protocolo común de integración
y sondeo de cuenca.


\section{Conclusiones}

El balance armónico y la continuación localizaron candidatos en Chua,
Kalman--Fitts y MAVPD; el PLL requirió un mapa de retorno sobre el cilindro.
Las obstrucciones de ganancia en Kalman--Fitts y de fase en el PLL determinaron
las rutas alternativas. La búsqueda por recurrencias localizó además el
candidato de Chua a partir de un muestreo global independiente.

Los diagnósticos distinguen el ciclo estable de MAVPD en \(\xi=3.1\) del
candidato caótico obtenido por continuación hacia
\(\xi=\MAVPDChaosXi\). Este último combina un exponente máximo positivo,
retornos irregulares y controles entre integradores, con las limitaciones
de tiempo finito descritas en \cref{int:subsec:mavpd-chaos}.
\ifHAFOBasinMaterial
Los sondeos comunes coincidieron entre cálculos y no alcanzaron las nubes
candidatas en los casos positivos. MAVPD en \(\xi=3.5\) proporcionó el
control negativo, con contactos desde vecindades de los equilibrios.
\fi

La cobertura incluye sistemas escalares D0, transformaciones T1 y
alternativas A1. Los modelos multicanal requieren otra formulación; para los
candidatos enteros pendientes faltan especificaciones del modelo o sondeos
de cuenca, según la \cref{int:tab:catalog}.
